\chapter{Calling Java}\label{app-java}

What calling C is for ga68 (Appendix~\ref{app-nest}), calling Java is
for a68j. Since a68j compiles to the Java virtual machine, an Algol 68
program can use the classes of the Java platform directly: the static
methods and fields of a class, and its objects with their constructors,
methods and fields. This puts the large library of the JDK, and any
library or class written in Java, at the disposal of the program. The
program runs wherever a Java virtual machine runs, as other Java
programs do. The facility is an extension of a68j and is enabled with the
option \code|-fnest-java|; a program that uses it runs only with a68j.
Without the option, a68j reports what ga68 reports for such a program:
\texttt{Java is not a valid language indication}.

\section*{Formal holes}

A Java method or field is named by a formal hole, written as for C,
with the language \code|Java| and the full name of the class:

\example{appH/methods}

\noindent The program is run with the option:
\begin{lstlisting}[style=shell]
$ java -jar a68j-backend.jar -fnest-java -run methods.a68
\end{lstlisting}
As for C, the type is not written in the hole but given by the context,
and a formal hole may only be written in a strong context. A hole of a
procedure type names a static method; a hole of another type, such as
\code|math_pi|, reads a static field once, when the declaration is
executed.

The difference from C is that Java classes describe themselves. a68j
looks the method up when it compiles the program and checks that the
declared type fits it. Here \code|cbrt| is declared with a \code|real|,
a 32-bit number, and \texttt{Math.cbrt} takes and returns a
\texttt{double}: a68j converts the values, where the C declaration of
Appendix~\ref{app-nest} would give wrong results. A declaration that no
method fits is an error when the program is compiled:

\example{appH/nosuch}

\noindent If a class has several methods of the name that fit, the
overloads of Java, a68j takes the one that needs the least conversion,
and reports an error if two are equally good. The parameter types can
also be written after the name of the method, as in
\code|nest Java "java.lang.String.join(CharSequence, Iterable)"|.

\section*{Types}

The types of Algol 68 and Java correspond as follows:

\begin{center}
\small
\begin{tabular}{ll}
\toprule
Algol 68 & Java \\
\midrule
\code|void| (as a result) & any result, which is discarded \\
\code|bool| & \texttt{boolean} \\
\code|char| & \texttt{int} (the character code) or \texttt{char} \\
\code|int|, \code|long int| & \texttt{int}, \texttt{long} \\
\code|bits| & \texttt{int} \\
\code|real|, \code|long real| & \texttt{float} or \texttt{double}, \texttt{double} \\
\code|string| & \texttt{String} \\
an array of these types & a Java array, copied \\
\code|ref| to an array & a Java array, copied back after the call \\
\code|proc| & an interface with one abstract method \\
\code|ref| to a record with one field & a Java object (see below) \\
\bottomrule
\end{tabular}
\end{center}

\noindent Where Java expects an object of a class such as \texttt{Object}
or \texttt{Integer}, a value of a simple type is passed as the
corresponding object, and an object returned by Java is converted back.

\section*{Functions as parameters}

Many Java methods take a function: a comparison for sorting, an action
for each element, a function to integrate. In Java such a parameter has
the type of an interface with a single abstract method, such as
\texttt{Comparator} or \texttt{DoubleUnaryOperator}. An Algol 68 function
of the matching type can be passed for it. \texttt{Arrays.sort} sorts the
array \code|names| with a routine text that compares the lengths of two
strings; since the parameter is a reference to the array, the sorted
array is copied back:

\example{appH/sort}

\section*{Java code of one's own}

Classes of one's own are found on the class path given with the option
\code|-fjava-classpath|. The class \texttt{Events} below keeps the
functions it receives and calls them later, and it has a method that
calls two functions while it runs: Newton's method for a zero of a
function \code|f| with the derivative \code|df|. The interface
\texttt{Handler} is one of its own; a function of type
\code|proc (string, int) void| fits it:

\javaprogram{appH/Events}

\noindent and the Algol 68 program that uses it:

\example{appH/callbacks}

\noindent The Java class is compiled with the Java compiler first:
\begin{lstlisting}[style=shell]
$ javac -d classes Events.java
$ java -jar a68j-backend.jar -fnest-java -fjava-classpath=classes -run callbacks.a68
\end{lstlisting}
The second handler changes the variable \code|total| of the program each
time \texttt{fire} calls it. Java receives Algol 68 functions only as
values that the program passes to it; it has no way to find a function of
the program by its name.

A function that Java keeps may be called after the block in which it was
declared has ended. In a68j this works, since its blocks keep their names
on the heap rather than on the stack, but it goes beyond the lifetime
rule of Section~\ref{sec-proc-lifetimes}, and a68j does not check it. The
functions that ga68 passes to C are bound to the stack frame of their
block (Appendix~\ref{app-nest}).

\section*{Objects}

A Java object is represented by a reference to a record with one field.
Each Java class gets a type of its own, declared once:

\example{appH/objects}

\noindent The record and its field are never used; they give the type
\code|List| a name of its own, so that a \code|List| cannot be confused
with a reference of another type. The value \code|nil| is Java's
\texttt{null}, and \code|is| compares two objects as \texttt{==} does in
Java. A name with \code|::| refers to a constructor (\code|::new|), a
method of an object, or a field of an object; the object is the first
parameter. The type arguments in angle brackets say what the type
parameters of a generic class stand for, as in Java: with
\code|ArrayList<String>|, a declaration of \code|add| with an \code|int|
parameter is an error when the program is compiled. A Java method can
itself be the function that is passed: \code|compare| is the method
\texttt{compareTo} of a \texttt{String}, whose object is its first
parameter, and \code|sort| receives it as the comparison.

\section*{Exceptions}

An exception that a Java method throws and does not catch ends the
program with a run-time error, at the position of the call in the Algol
68 program:

\example{appH/exception}

\noindent The program stops with the exit status 134, as for the
run-time errors of Algol 68 itself.

\section*{Programs compiled separately}

With \code|-d|, a68j writes the compiled program to a directory instead
of running it. The program is then run with the Java command, with the
runtime of a68j and the classes it uses on the class path (separated by
\texttt{;} on Windows):
\begin{lstlisting}[style=shell]
$ java -jar a68j-backend.jar -fnest-java -fjava-classpath=classes callbacks.a68 -d out
$ java -cp a68j-backend.jar:classes:out Callbacks
\end{lstlisting}
